\documentclass[12p]{ctexart}

\usepackage[a4paper,margin=0.8cm]{geometry}

\usepackage{tikz}
\usetikzlibrary{arrows.meta, positioning, fit, backgrounds}

\begin{document}

\pagestyle{empty}

\begin{center}
    {\LARGE \bfseries Cortex-A V.S. x86 架构的中断路径示意图}
\end{center}

\vspace{1em}

\begin{tikzpicture}[
    node distance=6mm and 18mm,
    >=Latex,
    block/.style={
        draw,
        rounded corners,
        align=center,
        minimum width=4.2cm,
        minimum height=9mm,
        font=\small
    },
    titlebox/.style={
        draw,
        rounded corners,
        fill=gray!10,
        minimum width=4.6cm,
        minimum height=10mm,
        font=\bfseries
    },
    groupbox/.style={
        draw,
        rounded corners,
        inner sep=8pt
    },
    note/.style={
        font=\small,
        align=left
    }
]

%--------------------------------
% Left column: Cortex-A
%--------------------------------
\node[titlebox] (armtitle) {Cortex-A / ARM 路径};

\node[block, below=of armtitle] (armdev) {外设 / 中断源\\Peripheral / IRQ Source};
\node[block, below=of armdev] (gicd) {GIC Distributor};
\node[block, below=of gicd] (gicc) {GIC CPU Interface\\或 Redistributor};
\node[block, below=of gicc] (armcore) {ARM Core};
\node[block, below=of armcore] (armentry) {异常入口\\Exception Entry};
\node[block, below=of armentry] (armirq) {Linux IRQ Handler};
\node[block, below=of armirq] (armthread) {IRQ Thread / RT Task\\（PREEMPT\_RT 可选）};

\draw[->, thick] (armdev) -- (gicd);
\draw[->, thick] (gicd) -- (gicc);
\draw[->, thick] (gicc) -- (armcore);
\draw[->, thick] (armcore) -- (armentry);
\draw[->, thick] (armentry) -- (armirq);
\draw[->, thick] (armirq) -- (armthread);

\node[note, left=8mm of armirq, anchor=east] (armnote) {
特点：\\
1. 路径相对简洁\\
2. GIC 负责中断分发\\
3. 实时性通常更稳定
};

\node[groupbox, fit=(armtitle)(armdev)(gicd)(gicc)(armcore)(armentry)(armirq)(armthread)] (armgroup) {};

%--------------------------------
% Right column: x86
%--------------------------------
\node[titlebox, right=1.5cm of armtitle] (x86title) {x86 路径};

\node[block, below=of x86title] (x86dev) {外设 / 中断源\\Peripheral / IRQ Source};
\node[block, below=of x86dev] (ioapic) {IO-APIC\\或 MSI / MSI-X};
\node[block, below=of ioapic] (lapic) {Local APIC};
\node[block, below=of lapic] (x86core) {x86 Core};
\node[block, below=of x86core] (idtentry) {IDT / Interrupt Gate\\中断入口};
\node[block, below=of idtentry] (x86irq) {Linux IRQ Handler};
\node[block, below=of x86irq] (x86thread) {IRQ Thread / RT Task\\（PREEMPT\_RT 可选）};

\draw[->, thick] (x86dev) -- (ioapic);
\draw[->, thick] (ioapic) -- (lapic);
\draw[->, thick] (lapic) -- (x86core);
\draw[->, thick] (x86core) -- (idtentry);
\draw[->, thick] (idtentry) -- (x86irq);
\draw[->, thick] (x86irq) -- (x86thread);

\node[note, right=8mm of x86irq, anchor=west] (x86note) {
特点：\\
1. APIC / MSI 体系成熟\\
2. 平均延迟可很低\\
3. 平台噪声因素更多\\
   （如 SMI / 电源管理）
};

\node[groupbox, fit=(x86title)(x86dev)(ioapic)(lapic)(x86core)(idtentry)(x86irq)(x86thread)] (x86group) {};

%--------------------------------
% Bottom summary
%--------------------------------
\node[note, below=16mm of armthread] (summary) {
% , anchor=north west
\textbf{说明：}\\
该图强调的是“中断从外设产生到 Linux 中断处理路径”的结构差异。\\
若使用 PREEMPT\_RT，Linux IRQ Handler 后面往往还会进入 IRQ 线程或唤醒实时任务。\\
Cortex-A 通常更强调确定性；x86 往往平均性能很强，但最坏延迟更需平台调优。
};

\end{tikzpicture}

\end{document}
